Un astronauta que és a l'espai vol mesurar la seva massa. Per fer-ho, s'asseu i es lliga a una cadira de 2,00~kg de massa que està unida a una molla de constant elàstica $k = 320\ \mathrm{N\,m^{-1}}$. L'astronauta s'impulsa i triga 62,8~s a fer 20 oscil·lacions completes.

\begin{apartat}{1}
Quina és la massa de l'astronauta?
\end{apartat}

\begin{apartat}{1}
Posteriorment aquest astronauta arriba a la Lluna, on fa oscil·lar un pèndol simple d'1,00~kg de massa i 1,50~m de longitud. Aquest pèndol triga 2~min i 1~s a fer 20 oscil·lacions completes. Quina és la intensitat del camp gravitatori a la superfície de la Lluna? Quina és la massa de la Lluna?
\end{apartat}

\dades{%
  $G = 6,67\times 10^{-11}\ \mathrm{N\,m^2\,kg^{-2}}$.\\
  $R_\mathrm{Lluna} = 1,737\times 10^{6}\ \mathrm{m}$.\\[2pt]
  El període d'oscil·lació d'un pèndol de longitud $L$ és $T = 2\pi\sqrt{\dfrac{L}{g}}$.}
